\documentclass[aspectratio=169,10pt]{beamer}

% =========================================================
% DEFENSE VERSION — revised for 15–20 minutes
% Silvia Busquets Xapellí — October 19, 2026
% =========================================================

\usepackage[utf8]{inputenc}
\usepackage[T1]{fontenc}
\usepackage{lmodern}
\usepackage{graphicx}
\usepackage{booktabs}
\usepackage{tabularx}
\usepackage{array}
\usepackage{amsmath}
\usepackage{xcolor}
\usepackage{hyperref}
\usepackage{ragged2e}

% ---------- Visual style ----------
\definecolor{FormalNavy}{HTML}{17324D}
\definecolor{AccentBlue}{HTML}{2F6F9F}
\definecolor{AccentGold}{HTML}{B98A2E}
\definecolor{SoftGrey}{HTML}{F4F6F8}
\definecolor{TextGrey}{HTML}{333333}

\usetheme{Boadilla}
\usecolortheme[named=FormalNavy]{structure}
\setbeamertemplate{navigation symbols}{}
\setbeamertemplate{footline}[frame number]
\setbeamertemplate{itemize item}{\color{FormalNavy}$\blacktriangleright$}
\setbeamertemplate{itemize subitem}{\color{AccentBlue}$\bullet$}
\setbeamercolor{title}{fg=FormalNavy}
\setbeamercolor{frametitle}{fg=FormalNavy}
\setbeamercolor{normal text}{fg=TextGrey}
\setbeamercolor{block title}{fg=white,bg=FormalNavy}
\setbeamercolor{block body}{fg=TextGrey,bg=SoftGrey}
\setbeamercolor{alerted text}{fg=AccentGold}
\setbeamerfont{title}{series=\bfseries,size=\Large}
\setbeamerfont{frametitle}{series=\bfseries,size=\large}

% ---------- Utility ----------
\newcommand{\bigclaim}[1]{%
  \begin{beamercolorbox}[wd=\linewidth,sep=0.7em,rounded=true]{block title}%
  \color{white}\large\bfseries #1
  \end{beamercolorbox}}

\newcommand{\metric}[2]{%
  \begin{minipage}{0.30\linewidth}\centering
  {\Large\bfseries\color{FormalNavy} #1}\\[-0.15em]
  {\scriptsize #2}
  \end{minipage}}

\newcommand{\safeimage}[2]{%
  \IfFileExists{#1}{\includegraphics[#2]{#1}}{%
    \fbox{\begin{minipage}[c][0.35\textheight][c]{0.78\linewidth}\centering
    Missing figure
    \end{minipage}}}}

\title[Birth Order and Academic Productivity]{Birth Order and Academic Productivity in Pre-Modern Europe}
\subtitle{Evidence from Brothers in Universities and Learned Academies}
\author{Silvia Busquets Xapellí}
\institute{Católica Lisbon School of Business \& Economics \;|\; UCLouvain\\
Supervisors: Prof. Hugo Reis and Prof. David de la Croix}
\date{Dissertation Defense \;|\; 19 October 2026}

\begin{document}

% =========================================================
% MAIN DECK
% =========================================================

% ---------- 1 ----------
\begin{frame}[plain]
  \titlepage
  \vspace{-1.0em}
  \begin{center}
  \footnotesize Remote dissertation defense
  \end{center}
\end{frame}

% ---------- 2 ----------
\begin{frame}{Research question and contribution}
\bigclaim{Did older brothers among pre-modern European scholars become more productive than their younger brothers?}
\vspace{0.6em}

\begin{columns}[T,onlytextwidth]
\begin{column}{0.48\textwidth}
\textbf{Why this is interesting}
\begin{itemize}
    \item Modern evidence often finds firstborn advantages in education and long-run outcomes.
    \item Historical evidence is more mixed and institutional context differs sharply from today.
    \item Scholarly families offer a setting in which long-run intellectual output can be observed.
\end{itemize}
\end{column}

\begin{column}{0.48\textwidth}
\textbf{What this dissertation adds}
\begin{itemize}
    \item A new sample of \textbf{645 rankable brother pairs}.
    \item Direct \textbf{within-family older -- younger comparisons}.
    \item A distinction between \textbf{entry into publication} and productivity conditional on publishing.
    \item Links to \textbf{academic institutions} and \textbf{regional inheritance regimes}.
\end{itemize}
\end{column}
\end{columns}
\end{frame}

% ---------- 3 ----------
\begin{frame}{Literature and hypotheses}
\small
\begin{columns}[T,onlytextwidth]
\begin{column}{0.48\textwidth}
\begin{block}{1. Birth order and human capital}
Modern within-family studies generally find an earlier-born advantage.\\
{\scriptsize Devereux et al. (2005); Behrman \& Taubman (1986); Kantarevic \& Mechoulan (2006)}
\end{block}

\begin{block}{2. Mechanisms}
Resource dilution, differential parental investment, and sibling competition.\\
{\scriptsize Becker \& Tomes (1976); Blake (1981); Price (2008); Hotz \& Pantano (2015)}
\end{block}
\end{column}

\begin{column}{0.48\textwidth}
\begin{block}{3. Historical and elite settings}
Historical effects are smaller or more context-dependent; elite formation and institutions matter.\\
{\scriptsize Clark \& Cummins (2024); Cools et al. (2024); De la Croix \& Goñi (2024)}
\end{block}

\begin{block}{4. Family institutions}
Inheritance systems can alter the allocation of resources within families.\\
{\scriptsize Todd (1990); Gay et al. (2024); Gobbi et al. (2026)}
\end{block}
\end{column}
\end{columns}

\vspace{0.25em}
\textbf{Ex ante hypothesis:} older-brother advantages should be larger where family resources are less equally distributed; institutional environments may also amplify sibling divergence.
\end{frame}

% ---------- 4 ----------
\begin{frame}{Data and sample construction}

\begin{columns}[T,onlytextwidth]

% ====================================================
% LEFT COLUMN — SOURCES
% ====================================================
\begin{column}{0.58\textwidth}

\textbf{Sources}

\vspace{0.15em}

\begin{itemize}
    \setlength{\itemsep}{0.30em}

    \item \textbf{RETE (Repertorium Eruditorum Totius Europae)}\\[-0.10em]
    {\small
    Scholars in European universities and learned academies;
    birth information and academic careers.}

    \item \textbf{Biographical repertories}\\[-0.10em]
    {\small
    Historical catalogues of professors and academic elites;
    additional family links (fathers, sons and brothers).}

    \item \textbf{VIAF (Virtual International Authority File)}\\[-0.10em]
    {\small
    Bibliographic productivity: titles, publishers,
    publication countries and coauthors.}

    \item \textbf{Wikipedia / Wikidata}\\[-0.10em]
    {\small
    Scholarly visibility (page length and languages) and
    academic affiliations / occupations.}

    \item \textbf{Todd (1990) + Gay et al. (2024)}\\[-0.10em]
    {\small
    Regional inheritance systems.}

\end{itemize}

\end{column}

% ====================================================
% RIGHT COLUMN — SAMPLE CONSTRUCTION
% ====================================================
\begin{column}{0.38\textwidth}

\textbf{Sample construction}

\vspace{0.10em}

\begin{center}

\begin{tabular}{ccc}
{\large\bfseries\color{FormalNavy}913}
&
{\large\bfseries\color{FormalNavy}645}
&
{\large\bfseries\color{FormalNavy}268}
\\[-0.10em]

{\tiny identified brother pairs}
&
{\tiny rankable pairs}
&
{\tiny unrankable pairs}

\end{tabular}

\end{center}

\vspace{0.15em}

\begin{block}{Rankable = usable for birth-order analysis}
\scriptsize
Both birth years are observed and different, so older and younger
brothers can be assigned without arbitrary choices.
\end{block}

\vspace{-0.25em}

\begin{alertblock}{Selection}
\scriptsize
Rankable pairs are more visible in the historical record than
unrankable pairs. Results therefore describe
\textbf{observed scholarly elites}, not all families.
\end{alertblock}

\end{column}

\end{columns}

\end{frame}

% ---------- 5 ----------
\begin{frame}{Outcome and transformation}

\begin{columns}[T,onlytextwidth]

% ====================================================
% LEFT COLUMN — OUTCOME + DISTRIBUTION
% ====================================================
\begin{column}{0.54\textwidth}

\begin{block}{Main outcome}
\small
\begin{itemize}
    \setlength{\itemsep}{0.20em}

    \item Number of titles attributed to each scholar in VIAF.

    \item Alternative outcomes capture visibility, publishers,
    coauthors, languages, and institutional involvement.
\end{itemize}
\end{block}

\vspace{0.10em}

\centering

\includegraphics[
    width=0.90\linewidth,
    height=0.34\textheight,
    keepaspectratio
]{figures/figure2_distribution_clean3.png}

\vspace{0.05em}

{\footnotesize
\textbf{\color{FormalNavy}Why transform?}
Titles are strongly right-skewed, with many observations close to zero.
}

\end{column}

% ====================================================
% RIGHT COLUMN — TRANSFORMATION
% ====================================================
\begin{column}{0.42\textwidth}

\begin{center}
{\footnotesize\textbf{Inverse hyperbolic sine transformation}}

\vspace{0.08em}

\[
\operatorname{asinh}(x)
=
\ln\!\left(x+\sqrt{x^2+1}\right)
\]
\end{center}

\vspace{0.05em}

\begin{block}{Key property}
\footnotesize
Defined at zero, while behaving approximately like a logarithm
for large positive values.
\end{block}

\vspace{0.10em}

\begin{alertblock}{Interpretation}
\footnotesize
Asinh coefficients are \textbf{not constant percentage effects near zero}.
Percentage interpretations depend on the chosen baseline.
\end{alertblock}

\end{column}

\end{columns}

\end{frame}

% ---------- 6 ----------
\begin{frame}{Descriptive within-family evidence}

\begin{columns}[T,onlytextwidth]

% ====================================================
% LEFT COLUMN — PAIRED COMPARISONS
% ====================================================
\begin{column}{0.57\textwidth}

\begin{block}{Paired comparison}
\footnotesize

For each brother pair:
\[
d_i = Outcome^{Older}_i - Outcome^{Younger}_i
\]

\[
H_0: E[d_i]=0
\]

\centering
\(E[d_i]>0\) $\Rightarrow$ older brothers have higher values on average.
\end{block}

\vspace{0.15em}

\centering

{\scriptsize
\setlength{\tabcolsep}{5pt}
\renewcommand{\arraystretch}{1.13}

\begin{tabular}{lrr}
\toprule
\textbf{Measure} & \textbf{Mean diff.} & \textbf{p-val.}\\
\midrule
Titles & 2.096 & .0076\\
Titles asinh & 0.245 & .0007\\
Has title & 0.056 & .0089\\
Wiki languages & 0.659 & .0215\\
Publishers & 0.667 & .0043\\
Coauthors & 1.056 & .0001\\
\bottomrule
\end{tabular}
}

\vspace{0.20em}

\begin{block}{Takeaway}
\footnotesize
Average differences generally favor older brothers,
but the advantage is \textbf{modest and far from universal}.
\end{block}

\end{column}

% ====================================================
% RIGHT COLUMN — HETEROGENEITY
% ====================================================
\begin{column}{0.39\textwidth}

\begin{center}
\textbf{Who is more productive?}

\vspace{0.20em}

\includegraphics[
    width=0.92\linewidth,
    height=0.36\textheight,
    keepaspectratio
]{figures/figure1_productivity_comparison.png}
\end{center}

\vspace{0.10em}

\begin{block}{Heterogeneity across families}
\footnotesize
\begin{tabular}{@{}lr@{}}
Older brother & \textbf{41.6\%}\\
Younger brother & \textbf{31.3\%}\\
Equal productivity & \textbf{27.1\%}
\end{tabular}
\end{block}

\vspace{0.15em}

{\footnotesize
The average advantage favors older brothers, but nearly
\textbf{6 in 10 pairs} do not show an older-brother advantage.
}

\end{column}

\end{columns}

\end{frame}

% ---------- 7 ----------
\begin{frame}{Where does the difference come from? Extensive vs.\ intensive margin}

\begin{columns}[T,onlytextwidth]

% ====================================================
% LEFT — EXTENSIVE
% ====================================================
\begin{column}{0.48\textwidth}

\textbf{Extensive margin: any recorded title}

\vspace{0.15em}

\[
Pr(Title_i>0)
=
\Lambda\!\left(
\alpha
+\beta Older_i
+\gamma AgeGap_i


+\delta FatherAcademia_i
\right))
\]

\begin{block}{Logit result}
\footnotesize
Older-sibling coefficient: \textbf{0.238***}\\
Odds ratio: \(e^{0.238}=\textbf{1.27}\)

\vspace{0.20em}

\textbf{Interpretation:}
older brothers have about \textbf{27\% higher odds}
of having at least one recorded title,
conditional on included controls.
\end{block}

{\scriptsize
\(N=1{,}290\) scholar observations
\(=645\) pairs \(\times 2\) brothers.\\
All scholars remain in the sample because the outcome is
whether \(Title>0\).
}

\end{column}

% ====================================================
% RIGHT — INTENSIVE
% ====================================================
\begin{column}{0.48\textwidth}

\textbf{Intensive margin: output among publishers}

\vspace{0.15em}

\[
asinh(Titles_i)
=
\alpha
+\beta Older_i
+\gamma AgeGap_i


+\delta FatherAcademia_i
+\varepsilon_i
\]

\begin{block}{OLS on positive-title sample}
\footnotesize
Older-sibling coefficient: \textbf{0.141}\\
Not statistically significant.
\end{block}

{\scriptsize
\(N=772\) scholars with positive title counts.\\
Scholars with zero titles are excluded because this specification
measures productivity \textbf{conditional on publishing}.
}

\vspace{0.35em}

\begin{alertblock}{Interpretation}
\footnotesize
The strongest birth-order pattern concerns
\textbf{participation in recorded scholarly production},
not output once a scholar is already publishing.
\end{alertblock}

\end{column}

\end{columns}

\end{frame}

% ---------- 8 ----------
\begin{frame}{Institutional environment and inheritance systems}

\setlength{\columnsep}{0.6em}

\begin{columns}[T,onlytextwidth]

% ====================================================
% LEFT — LARGE MAP + LEGENDS
% ====================================================
\begin{column}{0.70\textwidth}

\centering

\includegraphics[
    width=\linewidth,
    height=0.68\textheight,
    keepaspectratio
]{figures/figure3_map_tight.png}

\vspace{-0.10em}

\begin{columns}[T,onlytextwidth]

    \begin{column}{0.48\textwidth}
    \centering
    \includegraphics[
        width=0.95\linewidth,
        height=0.105\textheight,
        keepaspectratio
    ]{figures/figure3_legend_scholars_final.png}
    \end{column}

    \begin{column}{0.48\textwidth}
    \centering
    \includegraphics[
        width=0.95\linewidth,
        height=0.105\textheight,
        keepaspectratio
    ]{figures/figure3_legend_inheritance_final.png}
    \end{column}

\end{columns}

\end{column}

% ====================================================
% RIGHT — DEFINITIONS
% ====================================================
\begin{column}{0.28\textwidth}

\textbf{Inheritance}

\begin{itemize}
    \setlength{\itemsep}{0.10em}
    \setlength{\leftmargini}{1.15em}

    \item Assigned by birthplace region at NUTS3 level.
    \item Egalitarian: \textbf{30.1\%}
    \item Non-egalitarian: \textbf{46.5\%}
    \item Unknown: \textbf{23.4\%}
\end{itemize}

\vspace{0.08em}

\textbf{Institutions}

\begin{itemize}
    \setlength{\itemsep}{0.10em}
    \setlength{\leftmargini}{1.15em}

    \item Indicator \(=1\) if brothers are associated with
    different academic institutions.
\end{itemize}

\vspace{0.08em}

\begin{alertblock}{Important}
\scriptsize

``Older observed scholar'' is not necessarily the biological firstborn.

\vspace{0.08em}

The inheritance variable is a \textbf{regional proxy},
not a household rule.
\end{alertblock}

\end{column}

\end{columns}

\end{frame}

% ---------- 9 ----------
\begin{frame}{Empirical specification}
\small
\[
\Delta P_i =
\operatorname{asinh}(Titles^{Older}_i)-
\operatorname{asinh}(Titles^{Younger}_i)
\]
\vspace{-0.3em}
\[
\begin{aligned}
\Delta P_i ={}& \beta_0+\beta_1 AgeGap_i+\beta_2 FatherAcademia_i
+\beta_3 DifferentInstitution_i\\
&+\beta_4 NonEgalitarian_i+\beta_5 Unknown_i
+\sum_{c=13}^{17}\gamma_c Century_{ic}+\varepsilon_i
\end{aligned}
\]
\vspace{-0.5em}
\begin{itemize}\setlength\itemsep{0.25em}
  \item Positive \(\Delta P_i\): higher recorded productivity for the older brother.
  \item OLS with standard errors clustered at family level; 18th century omitted.
  \item Pair differencing removes shared family-level levels, but does not randomize birth order, survival, observation, career choice or institution.
\end{itemize}
\vspace{-0.2em}
\begin{alertblock}{Identification claim}
I interpret these as \textbf{within-family associations}, not fully causal estimates.
\end{alertblock}
\end{frame}

% ---------- 10 ----------
\begin{frame}{Preferred regression: institutions are the most robust correlate}
\begin{table}
\centering
\small
\begin{tabular}{lccc}
\toprule
Variable & Coefficient & S.E. & Significance\\
\midrule
Age gap & 0.019 & 0.011 & *\\
Father in academia & -0.163 & 0.150 & n.s.\\
Different institution & \textbf{0.485} & 0.180 & \textbf{***}\\
Non-egalitarian inheritance & 0.135 & 0.175 & n.s.\\
Unknown inheritance & 0.130 & 0.205 & n.s.\\
\bottomrule
\end{tabular}
\end{table}

\vspace{0.35em}
\begin{block}{What I would emphasize}
The coefficient on \textbf{different institution} is positive and statistically robust: brothers associated with different institutions display larger older -- younger productivity gaps.
\end{block}

\begin{alertblock}{What I would not claim}
The coefficient does not establish that moving brothers to different institutions \emph{causes} a 63\% increase in productivity. It is an association in transformed gap units.
\end{alertblock}

{\scriptsize Preferred model: \(N=645\), \(R^2=0.035\), century indicators included, family-clustered SE.}
\end{frame}

% ---------- 11 ----------
\begin{frame}{Historical heterogeneity}

\small

% ====================================================
% TOP — FIGURE + PERIOD COMPARISON
% ====================================================
\begin{columns}[T,onlytextwidth]

% ---------------- LEFT: FIGURE ----------------
\begin{column}{0.59\textwidth}

\centering

\safeimage{figures/historical_heterogeneity_clear.png}{
    width=\linewidth,
    height=0.50\textheight,
    keepaspectratio
}

\end{column}

% ---------------- RIGHT: RESULTS ----------------
\begin{column}{0.37\textwidth}

\begin{block}{Period comparison}
\scriptsize

\textbf{Before the 16th century}
\begin{itemize}
    \setlength{\itemsep}{0.05em}
    \item \textbf{33 pairs}
    \item Age gap: \textbf{0.110***}
    \item Different institution: 0.995, imprecise
\end{itemize}

\vspace{0.15em}

\textbf{16th century onward}
\begin{itemize}
    \setlength{\itemsep}{0.05em}
    \item \textbf{612 pairs}
    \item Age gap: 0.013, n.s.
    \item Different institution: \textbf{0.474**}
\end{itemize}

\end{block}

\end{column}

\end{columns}

\vspace{0.20em}

% ====================================================
% BOTTOM — MOTIVATION + CAUTION
% ====================================================
\begin{columns}[T,onlytextwidth]

% ---------------- WHY SPLIT ----------------
\begin{column}{0.59\textwidth}

\begin{alertblock}{Why split at the 16th century?}
\scriptsize
Inheritance-system variation is very limited in the earliest centuries:
\textbf{non-egalitarian regimes are almost absent in the 13th and 14th centuries}.

The later period therefore provides substantially more variation
in inheritance regimes.
\end{alertblock}

\end{column}

% ---------------- CAUTION ----------------
\begin{column}{0.37\textwidth}

\begin{alertblock}{Caution}
\scriptsize
The early-period estimates are based on only
\textbf{33 brother pairs}.

They are informative, but substantially less precise and should
therefore be interpreted cautiously.
\end{alertblock}

\end{column}

\end{columns}

\end{frame}

% ---------- 12 ----------
\begin{frame}{Robustness and main threats to interpretation}

\small

\begin{columns}[T,onlytextwidth]

% ====================================================
% LEFT — ROBUSTNESS
% ====================================================
\begin{column}{0.48\textwidth}

\textbf{Robustness: what does each test check?}

\vspace{0.10em}

\begin{itemize}
    \setlength{\itemsep}{0.22em}

    \item \textbf{Age gap $\leq 20$ years}\\[-0.08em]
    {\scriptsize Excludes unusually large sibling age gaps.}

    \item \textbf{Two observed sons}\\[-0.08em]
    {\scriptsize Avoids multiple overlapping brother pairs in larger families.}

    \item \textbf{Alternative outcomes}\\[-0.08em]
    {\scriptsize Checks that results are not specific to VIAF titles.}

    \item \textbf{Alternative functional forms}\\[-0.08em]
    {\scriptsize Tests dependence on the asinh transformation.}

    \item \textbf{Country controls}\\[-0.08em]
    {\scriptsize Checks whether geographic composition drives the result.}

\end{itemize}

\vspace{0.10em}

\begin{block}{Stable pattern}
\scriptsize
Institutional divergence coefficients remain positive;
inheritance coefficients remain weak.
\end{block}

\end{column}

% ====================================================
% RIGHT — LIMITATIONS
% ====================================================
\begin{column}{0.48\textwidth}

\textbf{Main threats / limitations}

\vspace{0.10em}

\begin{itemize}
    \setlength{\itemsep}{0.20em}

    \item \textbf{Historical selection}\\[-0.08em]
    {\scriptsize Results apply to observed scholarly elites.}

    \item \textbf{Rankability selection}\\[-0.08em]
    {\scriptsize Better-documented scholars are overrepresented.}

    \item \textbf{VIAF measurement}\\[-0.08em]
    {\scriptsize Titles also reflect preservation, editions and cataloguing.}

   \item \textbf{Endogenous institutions}\\[-0.08em]
    {\scriptsize Institution choice may reflect ability or opportunities.}

    \item \textbf{Inheritance proxy}\\[-0.08em]
    {\scriptsize Regional regimes may differ from household practice.}

\end{itemize}

\end{column}

\end{columns}

\end{frame}

% ---------- 13 ----------
\begin{frame}{Conclusion}
\begin{enumerate}
  \item There is a \textbf{modest older-brother advantage} in recorded scholarly productivity.
  \item The clearest birth-order pattern is on the \textbf{extensive margin}: older brothers have higher odds of having at least one recorded title.
  \item Conditional on publishing, the older-sibling coefficient is small and statistically insignificant.
  \item \textbf{Institutional divergence} is more consistently associated with within-family productivity gaps than regional inheritance regimes.
  \item The relationship is historically heterogeneous: age gaps matter more in the very small early sample; institutional divergence is more robust from the 16th century onward.
\end{enumerate}

\vspace{0.5em}
\bigclaim{The evidence points less to a universal primogeniture effect than to birth-order differences shaped by access to scholarly activity and institutional context.}
\end{frame}

% ---------- 14 ----------
\begin{frame}[plain]
\begin{center}
\vspace{2.0em}
{\Huge\color{FormalNavy}\textbf{Thank you}}\\[1.2em]
{\Large Questions and discussion}\\[1.8em]
{\small Silvia Busquets Xapellí}
\end{center}
\end{frame}

% =========================================================
% BACKUP SLIDES
% =========================================================
\appendix

% ---------- B1 ----------
\begin{frame}{Backup: full preferred regression}
\centering
\scriptsize
\begin{tabular}{lcc}
\toprule
Variable & Coefficient & S.E.\\
\midrule
Constant & -0.162 & 0.197\\
Age gap & 0.019* & 0.011\\
Father in academia & -0.163 & 0.150\\
Different institution & 0.485*** & 0.180\\
Non-egalitarian inheritance & 0.135 & 0.175\\
Unknown inheritance & 0.130 & 0.205\\
13th century & 0.099 & 0.223\\
14th century & 1.764*** & 0.552\\
15th century & 0.863** & 0.345\\
16th century & 0.025 & 0.226\\
17th century & 0.099 & 0.161\\
\midrule
Observations & \multicolumn{2}{c}{645}\\
\(R^2\) & \multicolumn{2}{c}{0.035}\\
Adj. \(R^2\) & \multicolumn{2}{c}{0.020}\\
\bottomrule
\end{tabular}

\vspace{0.4em}
{\scriptsize Dependent variable: \(\operatorname{asinh}(Titles_O)-\operatorname{asinh}(Titles_Y)\). 18th century omitted. Family-clustered SE.}
\end{frame}

% ---------- B2 ----------
\begin{frame}{Backup: sample selection}
\begin{columns}[T,onlytextwidth]
\begin{column}{0.52\textwidth}
\centering
\small
\begin{tabular}{lcc}
\toprule
 & Rankable & Unrankable\\
\midrule
Brother pairs & 645 & 268\\
Mean titles & 10.84 & 1.87\\
Individuals with no titles & 40.2\% & 81.6\%\\
Pairs with no titles & 25.4\% & 69.3\%\\
Father in academia & 38.8\% & 41.0\%\\
\bottomrule
\end{tabular}
\end{column}

\begin{column}{0.44\textwidth}
\begin{itemize}
  \item Birth order is observable disproportionately for better documented scholars.
  \item This can affect external validity.
  \item It is not solved by the within-family comparison.
\end{itemize}
\end{column}
\end{columns}
\end{frame}

% ---------- Backup: concrete heterogeneity examples ----------
\begin{frame}{Backup: concrete examples from the data}

\begin{columns}[T,onlytextwidth]

% ====================================================
% LEFT — SMALL GAP
% ====================================================
\begin{column}{0.47\textwidth}

\begin{block}{Bernoulli brothers: similar productivity}
\small

\textbf{Jacob Bernoulli}\\
Older brother: \textbf{63 VIAF titles}

\vspace{0.20em}

\textbf{Johann Bernoulli}\\
Younger brother: \textbf{59 VIAF titles}

\vspace{0.25em}

{\footnotesize
A case where birth order is associated with only a very small
within-family productivity difference.
}

\end{block}

\end{column}

% ====================================================
% RIGHT — LARGE GAP
% ====================================================
\begin{column}{0.47\textwidth}

\begin{alertblock}{Thomasius brothers: large productivity gap}
\small

\textbf{Christian Thomasius}\\
Older brother: \textbf{71 VIAF titles}

\vspace{0.20em}

\textbf{Gottfried Thomasius}\\
Younger brother: \textbf{5 VIAF titles}

\vspace{0.25em}

{\footnotesize
One of the largest within-family productivity gaps
observed in the dataset.
}

\end{alertblock}

\end{column}

\end{columns}

\vspace{0.30em}

\begin{block}{Schott brothers: low visibility, but observable in structured data}
\footnotesize

\textbf{August Ludwig Schott (1751--1787): 4 titles}
\qquad
\textbf{Andreas Heinrich Schott (1758--1831): 6 titles}

\vspace{0.12em}

Neither brother has a significant presence in accessible historical sources.
The pair is identified through \textbf{prosopographical data}, illustrating how structured
historical datasets recover scholars who would otherwise be difficult to observe.

\end{block}

\vspace{0.15em}

{\footnotesize
\textbf{\color{FormalNavy}Takeaway:}
the dataset contains both highly visible scholarly families and much less visible pairs,
with very different within-family productivity patterns.
}

\end{frame}




% ------------------------------------------------------------
% Backup: Scholarly Productivity Across Samples
% ------------------------------------------------------------

\begin{frame}{Backup: Scholarly Productivity Across Samples}

\small

\begin{table}
\centering
\begin{tabular}{lcc}
\hline
\textbf{Sample} &
\textbf{Mean Titles} &
\textbf{Observations} \\
\hline
All scholars in the dataset       & 3.81  & 91,931 \\
Scholars with identified siblings & 8.29  & 1,552 \\
Older brothers (rankable)         & 11.89 & 645 \\
Younger brothers (rankable)       & 9.79  & 645 \\
\hline
\end{tabular}
\end{table}

\vspace{0.35cm}

\begin{block}{Key Takeaways}
\begin{itemize}
    \item Scholars with identified siblings have higher
    recorded productivity than the overall sample.
    \item Older brothers have, on average, 2.10 more
    recorded titles than younger brothers.
    \item The median number of titles in the full
    dataset is zero.
\end{itemize}
\end{block}

\vspace{0.15cm}

{\scriptsize
\textit{Notes:} The full historical academic dataset
includes 91,931 scholars. Among the 1,556 unique
scholars with identified siblings, title data are
available for 1,552 (four have missing values).
Zero-title observations are included in the means.
Older and younger brother means are based on
645 rankable pairs rather than unique individuals.
}

\end{frame}


% ---------- B3 ----------
\begin{frame}{Backup: asinh interpretation}

\small

\begin{block}{Transformations}
\[
\operatorname{asinh}(x)
=
\ln\!\left(x+\sqrt{x^2+1}\right)
\]

\[
\sinh(z)
=
\frac{e^z-e^{-z}}{2}
\]
\end{block}

\vspace{0.15em}

\textbf{Percentage illustration}

\[
\%\Delta y =
100\left[
\frac{
\sinh\!\left(
\operatorname{asinh}(y_0)+\beta
\right)
}{
y_0
}
-1
\right]
\]

\begin{itemize}
    \setlength{\itemsep}{0.25em}

    \item This is a \textbf{baseline-dependent illustration},
    not a constant semi-elasticity.

    \item Using \(y_0=9.79\):

    \begin{itemize}
        \item \(\beta=0.245
        \Rightarrow
        \%\Delta y \approx \textbf{27.9\%}\)

        \item \(\beta=0.485
        \Rightarrow
        \%\Delta y \approx \textbf{62.7\%}\)
    \end{itemize}

\end{itemize}

\begin{alertblock}{Interpretation}
\footnotesize
The percentage changes depend on the chosen baseline \(y_0\).
For this reason, the main results should be reported first in
\textbf{asinh units}, with percentage conversions used only as illustrations.
\end{alertblock}

\end{frame}

% ---------- B4 ----------
\begin{frame}{Backup: extensive and intensive margins}

\centering
\small

\begin{tabular}{lcc}
\toprule
 & Extensive (Logit) & Intensive (OLS)\\
 & \(Pr(Title>0)\) & Asinh titles\\
\midrule
Constant & 0.051 & 2.892***\\
Older sibling & \textbf{0.238***} & 0.141\\
Age gap & -0.002 & -0.013\\
Father in academia & 0.654*** & 0.090\\
\midrule
Observations & 1290 & 772\\
Family-clustered SE & Yes & Yes\\
\bottomrule
\end{tabular}

\vspace{0.6em}

\begin{block}{Logit interpretation}
\footnotesize

\[
OR=e^{\beta}=e^{0.238}\approx1.27
\]

\[
(1.27-1)\times100\approx\textbf{27\% higher odds}
\]

So older brothers have about \textbf{27\% higher odds}
of having at least one recorded title, conditional on the included controls.

\end{block}

\begin{alertblock}{Important}
\footnotesize
This is an increase in \textbf{odds}, not a 27\% increase in probability.
\end{alertblock}

\end{frame}

% ---------- B5 ----------
\begin{frame}{Backup: period-specific regression}
\centering
\small
\begin{tabular}{lcc}
\toprule
 & Before 16th c. & 16th c. onward\\
\midrule
Observations & 33 & 612\\
Age gap & 0.110*** & 0.013\\
Father in academia & 0.415 & -0.160\\
Different institution & 0.995 & 0.474**\\
Non-egalitarian inheritance & -0.719 & 0.164\\
Unknown inheritance & -1.730 & 0.175\\
\(R^2\) & 0.387 & 0.020\\
Adj. \(R^2\) & 0.215 & 0.008\\
\bottomrule
\end{tabular}
\end{frame}

% ---------- B6 ----------
\begin{frame}{Backup: robustness across functional forms}
\small
\begin{tabularx}{\linewidth}{>{\bfseries}p{0.23\linewidth}X X}
\toprule
Specification & Different institution & Main caution\\
\midrule
Asinh gap & 0.485*** & Baseline-dependent percentage translation\\
Log gap, \(\epsilon=0.1\) & 0.577** & Depends on zero adjustment\\
Raw title gap & 5.160** & More sensitive to outliers\\
Age gap \(\leq20\) & 0.478*** & \(N=624\)\\
Two-son families & 0.493** & \(N=467\)\\
Country controls & 0.388* & Many countries have small cell sizes\\
\bottomrule
\end{tabularx}

\vspace{0.6em}
\begin{block}{Takeaway}
The sign and broad qualitative pattern for institutional divergence survive several alternative specifications, although precision and magnitude vary.
\end{block}
\end{frame}

% ---------- B7 ----------
\begin{frame}{Backup: literature map}
\begin{itemize}\setlength\itemsep{0.35em}
\item \textbf{Modern birth order:} Devereux et al. (2005), Behrman \& Taubman (1986), Kantarevic \& Mechoulan (2006). Earlier-born children tend to have better schooling outcomes in modern within-family data.
\item \textbf{Mechanisms:} Becker \& Tomes (1976), Blake (1981), Price (2008), Hotz \& Pantano (2015). Finite resources, parental time, strategic responses and sibling competition.
\item \textbf{Historical evidence:} Clark \& Cummins (2024), Cools et al. (2024), Nuevo-Chiquero et al. (2022). Historical birth-order effects are mixed and context-dependent.
\item \textbf{Academic elites:} De la Croix \& Goñi (2024), De la Croix et al. (2024). Family background, institutions and networks matter in historical academia.
\item \textbf{Inheritance institutions:} Todd (1990), Gay et al. (2024), Gobbi et al. (2026). Regional family systems vary spatially and can affect household incentives.
\end{itemize}
\end{frame}

% ---------- B8 ----------
\begin{frame}{Backup: likely econometrics questions}
\begin{itemize}\setlength\itemsep{0.35em}
\item \textbf{Causal?} Within-family differencing removes shared family levels, but not selection, sibling-specific shocks, endogenous institution choice or measurement error.
\item \textbf{Why cluster by family?} Larger families generate several pair observations; family clustering allows residual correlation within reconstructed families.
\item \textbf{Why not count models?} The estimate is an older--younger within-pair gap including zeros; asinh gives a symmetric transformed difference. Count models are a useful extension, not the same estimand.
\item \textbf{Why not pair fixed effects?} The regression is already at pair level; the dependent variable is the within-pair difference, so shared pair-level levels cancel.
\item \textbf{Low $R^2$?} The goal is not prediction; historical productivity is highly heterogeneous. The question is whether selected covariates are systematically associated with the gap.
\end{itemize}
\end{frame}

% ---------- B9 ----------
\begin{frame}{Backup: questions about measurement and selection}
\begin{itemize}\setlength\itemsep{0.35em}
\item \textbf{Does VIAF measure true productivity?} No. It is a bibliographic proxy for recorded output and also reflects editions, translations, preservation and cataloguing.
\item \textbf{Why compare scholars only?} The research question concerns observed academic elites; it does not estimate birth-order effects for the general population.
\item \textbf{Is older the same as firstborn?} No. I identify relative age among observed scholarly brothers, not complete biological birth rank.
\item \textbf{Could different institutions reflect ability?} Yes. Institution choice is endogenous; I interpret it as an association with divergent trajectories, not a causal institution effect.
\item \textbf{Could rankability bias results?} Yes. Rankable pairs have richer records, so the external validity is selected and should be stated explicitly.
\end{itemize}
\end{frame}


\end{document}
